\documentclass[aspectratio=169, 10pt]{beamer}

\usetheme{Madrid}
\usefonttheme{professionalfonts}

\usepackage[utf8]{inputenc}
\usepackage[indonesian]{babel}
\usepackage{amsmath,amsfonts,amssymb}
\usepackage{graphicx}
\usepackage{booktabs}
\usepackage{colortbl}
\usepackage{listings}
\usepackage{xcolor}

\definecolor{mainblue}{RGB}{24, 43, 73}
\definecolor{accentblue}{RGB}{41, 128, 185}
\definecolor{accentgold}{RGB}{230, 126, 34}
\definecolor{darkgray}{RGB}{44, 62, 80}
\definecolor{lightbg}{RGB}{248, 250, 252}
\definecolor{bordergray}{RGB}{200, 210, 220}
\definecolor{codebg}{RGB}{240, 243, 246}
\definecolor{codegreen}{RGB}{39, 174, 96}

\setbeamercolor{structure}{fg=mainblue}
\setbeamercolor{palette primary}{bg=mainblue, fg=white}
\setbeamercolor{palette secondary}{bg=accentblue, fg=white}
\setbeamercolor{palette tertiary}{bg=mainblue!80, fg=white}
\setbeamercolor{title}{fg=white, bg=mainblue}
\setbeamercolor{frametitle}{fg=white, bg=mainblue}
\setbeamercolor{block title}{bg=mainblue, fg=white}
\setbeamercolor{block body}{bg=lightbg, fg=darkgray}
\setbeamercolor{block title alerted}{bg=accentgold, fg=white}
\setbeamercolor{block body alerted}{bg=accentgold!10, fg=darkgray}

\lstset{
    language=Python,
    backgroundcolor=\color{codebg},
    basicstyle=\ttfamily\scriptsize\color{darkgray},
    keywordstyle=\bfseries\color{accentblue},
    stringstyle=\color{accentgold},
    commentstyle=\itshape\color{codegreen},
    showstringspaces=false,
    breaklines=true,
    frame=single,
    rulecolor=\color{bordergray},
    numbers=none
}

\title[\textbf{Deret Waktu (Minggu 13)}]{\textbf{Deteksi Anomali \& Kegagalan Sistem Tenaga}}
\subtitle{Isolation Forest, Rekonstruksi LSTM Autoencoder, dan Deteksi Pencurian Listrik}
\author[\textbf{Ir. Novalio Daratha, Ph.D.}]{\textbf{Ir. Novalio Daratha, S.T., M.Sc., Ph.D.}}
\institute[FT UNIB]{Program Studi Teknik Elektro -- Fakultas Teknik\\Universitas Bengkulu}
\date{Minggu 13 / Semester Ganjil 2026}

\begin{document}

\begin{frame}[plain]
    \titlepage
\end{frame}

\begin{frame}{Agenda Perkuliahan Minggu ke-13}
    \begin{columns}[T]
        \begin{column}{0.48\textwidth}
            \begin{block}{Capaian Pembelajaran (Sub-CPMK 12)}
                Mahasiswa mampu mengidentifikasi taksonomi anomali pada data sensor kelistrikan, merancang model deteksi tanpa pengawasan (*Unsupervised Anomaly Detection*) berbasis Isolation Forest dan Autoencoder, serta mengevaluasi metrik presisi/recall.
            \end{block}
        \end{column}
        \begin{column}{0.48\textwidth}
            \begin{enumerate}
                \item \textbf{Taksonomi Anomali Data Kelistrikan}
                \item \textbf{Metode Statistik: Z-Score \& CUSUM}
                \item \textbf{Algoritma Isolation Forest}
                \item \textbf{Deep Learning: LSTM Autoencoder}
                \item \textbf{Metrik Evaluasi: Precision, Recall, \& F1-Score}
                \item \textbf{Studi Kasus: Deteksi Pencurian Smart Meter}
            \end{enumerate}
        \end{column}
    \end{columns}
\end{frame}

\begin{frame}{Taksonomi Anomali pada Sinyal Sistem Tenaga}
    \begin{columns}[T]
        \begin{column}{0.33\textwidth}
            \begin{block}{1. Point Anomalies}
                Titik tunggal yang melompat ekstrem sesaat (misal sambaran petir, spike tegangan transien).
            \end{block}
        \end{column}
        \begin{column}{0.33\textwidth}
            \begin{block}{2. Contextual Anomalies}
                Nilai berada dalam batas normal namun aneh dalam konteks waktu (misal beban puncak 500 kW terjadi pada pukul 03.00 pagi).
            \end{block}
        \end{column}
        \begin{column}{0.33\textwidth}
            \begin{block}{3. Collective Anomalies}
                Sekumpulan urutan yang polanya menyimpang (misal penurunan flat konsumsi meter pintar akibat manipulasi bypass kWh).
            \end{block}
        \end{column}
    \end{columns}
\end{frame}

\begin{frame}{Deteksi Anomali Berbasis Rekonstruksi (LSTM Autoencoder)}
    Autoencoder dilatih untuk memadatkan dan merekonstruksi \textbf{hanya pola sinyal normal}:
    \begin{block}{Prinsip Galat Rekonstruksi}
        \begin{equation}
            \text{Error Rekonstruksi } \mathcal{L}(t) = \| \mathbf{x}_t - \hat{\mathbf{x}}_t \|^2
        \end{equation}
        \begin{itemize}
            \item \textbf{Kondisi Normal}: Jaringan mampu merekonstruksi dengan sangat akurat ($\mathcal{L}(t) < \tau$).
            \item \textbf{Kondisi Anomali / Rusak}: Jaringan gagal merekonstruksi pola yang belum pernah dipelajari, sehingga $\mathcal{L}(t) > \tau$ (Alarm Anomali!).
        \end{itemize}
    \end{block}
    \begin{center}
        \textbf{Penentuan Ambang Batas}: $\tau = \mu_{\text{train\_error}} + 3 \cdot \sigma_{\text{train\_error}}$.
    \end{center}
\end{frame}

\begin{frame}[fragile]{Implementasi Python: LSTM Autoencoder Anomaly Detector}
\begin{lstlisting}
import numpy as np
from sklearn.metrics import precision_recall_fscore_support

# Model Autoencoder dilatih pada data operasi normal (X_train_normal)
autoencoder.fit(X_train_normal, X_train_normal, epochs=50, batch_size=32, verbose=False)

# Evaluasi pada data uji (mengandung anomali pencurian daya)
X_pred = autoencoder.predict(X_test)
mse_loss = np.mean(np.power(X_test - X_pred, 2), axis=(1, 2))

# Ambang Batas Deteksi Dinamis
threshold = np.mean(train_loss) + 3.0 * np.std(train_loss)
y_pred_anomaly = (mse_loss > threshold).astype(int)

# Metrik Evaluasi
p, r, f1, _ = precision_recall_fscore_support(y_test_ground_truth, y_pred_anomaly, average='binary')
print(f"Hasil Deteksi -> Precision: {p:.3f}, Recall: {r:.3f}, F1-Score: {f1:.3f}")
\end{lstlisting}
\end{frame}

\begin{frame}{Rangkuman \& Persiapan Minggu ke-14}
    \begin{block}{Poin Kunci Minggu 13}
        \begin{itemize}
            \item Anomali kelistrikan memiliki karakteristik kontekstual dan kolektif yang unik.
            \item Isolation Forest unggul untuk anomali titik dan deteksi cepat tanpa asumsi distribusi.
            \item LSTM Autoencoder merepresentasikan standar terbaik untuk mendeteksi manipulasi pencurian listrik dan degradasi kesehatan aset transformator.
        \end{itemize}
    \end{block}
    \vspace{0.2cm}
    \begin{alertblock}{Materi Minggu Depan (Minggu ke-14)}
        \textbf{Integrasi Peramalan pada Operasi Sistem Tenaga}: Unit Commitment, penjadwalan baterai BESS pada smart microgrid, dan kuantisasi interval ketidakpastian.
    \end{alertblock}
\end{frame}

\begin{frame}[plain]
    \begin{center}
        \vspace{1.5cm}
        {\Huge \color{mainblue} \textbf{Terima Kasih}}\\[0.4cm]
        {\large Pertanyaan \& Diskusi Deteksi Anomali Sistem Kelistrikan}\\[1.2cm]
        \textbf{Laboratorium Sistem Tenaga \& Komputasi Cerdas}\\
        Jurusan Teknik Elektro -- Universitas Bengkulu
    \end{center}
\end{frame}

\end{document}
